\chapter{Synthetic Nervous System Simulation}\label{ch:sns}
\graphicspath{{D:/Github/Bipedal_Robot/Documentation/Reports and Papers/Dissertation/Figures/60-sns/}{figs/60-sns/}{D:/Github/Bipedal_Robot/Documentation/Reports and Papers/Dissertation/Figures/60-sns/AnimatLab_reference/}{figs/60-sns/AnimatLab_reference/}}

The preceding chapters validate the robot's mechanics: a force model for the braided pneumatic actuators (Chapter~\ref{ch:muscles}) and a corrected model of the joint torque those actuators deliver through the biomimetic knee (Chapter~\ref{ch:torque}). This chapter builds the other half of the neuromechanical system, the controller, together with the simulation tools required to develop it. The controller is a synthetic nervous system (SNS): a network of conductance-based neuron and synapse models whose organization follows the mammalian spinal cord \citep{rybak_modelling_2006, mccrea_organization_2008, deng_neuromechanical_2019}. Its parameters are not trained in the machine-learning sense but set from physiological data and tuned against explicit verification tests, which keeps every neuron an identifiable cell class and lets the simulated system support lesion, stimulation, and sensory-deletion experiments that cannot be performed on living subjects (Section~\ref{sec:bg_ethics}). The neural vocabulary, the two-level central pattern generator, and the afferent channels are defined in Sections~\ref{sec:bg_neural}--\ref{sec:bg_sns} and used here without redefinition; the neuron and synapse equations themselves are recorded in Appendix~\ref{app:neuron_equations}.

The chapter tells one story in order. Section~\ref{sec:prior_walker} records the AnimatLab bipedal walker I built with Morrow at the start of this work, which combined two published designs and stepped in air but never walked. Section~\ref{sec:sns_gaps} states the four capability gaps that attempt exposed. Sections~\ref{sec:preliminary_sim_methods} and~\ref{sec:preliminary_sim_results} present the new software tools built to close those gaps and the verification tests that establish what the tools do and do not show, with the spiking-neuron variants of the same networks following in Section~\ref{sec:spiking_results}. Section~\ref{sec:sns_discussion} closes with the delivered capabilities, their limits, and the program they enable for Chapter~\ref{ch:futurework}.

\section{The AnimatLab Bipedal Walker}\label{sec:prior_walker}

As preliminary work of this dissertation, built before the simulation toolchain of this chapter existed, Morrow and I combined two published designs into a neuromechanical model of a bipedal walker in AnimatLab, a graphical environment for building biomechanical systems actuated by synthetic neural systems \citep{cofer_animatlab_2010}: the biomechanics came from the suspended bipedal walker of \citet{li_bipedal_2017}, and the controller from the two-layer CPG organization of \citet{deng_neuromechanical_2019}. That combination is the direct ancestor of the work reported here: its biomechanics and its controller architecture carry through every model in this chapter, and the limitations of its hand-tuned parameters are the reason the present toolchain emphasizes systematic verification.

The biomechanical model consisted of nine rigid bodies (pelvis, thigh, shank, foot, and toe segments) connected by sagittal-plane hinge joints with human-like rotational capabilities and actuated by antagonistic pairs of Hill-type muscles at the hip, knee, and ankle. The toe joint, my addition to the Li design, carried a passive spring instead of a muscle pair, simplifying the system. The Li model itself was driven by pattern generators entrained by heel and toe mechanosensory receptors and hip-angle position sensors standing in for the group-I loading, cutaneous, and hip-joint afferents of the biological scheme; muscle afferents were not modeled. It also required support against falling, because a system with only planar joints cannot actively balance side to side.

The synthetic nervous network was built from non-spiking, leaky-integrator neurons. A rhythm-generator layer set the timing of the system; a pattern-formation layer reshaped the timing signals into phase-appropriate activation of antagonist muscle groups; and a sensorimotor layer containing Renshaw-cell-like recurrent inhibition drove the motor neurons. The sensorimotor layer received Ia (length and velocity) and Ib (tension) afferent feedback from the muscles, but the feedback gains were never systematically tuned, and no ground-contact signals were modeled. Every applied current, synaptic weight, and muscle-tension parameter was set by hand: each was adjusted by trial and error, one at a time, until the three layers oscillated in phase and the motor neurons produced stepping-like activation of the muscle pairs. That hand-tuning consumed the bulk of the project's effort, left no record of the sensitivity of the behavior to individual parameters, and did not generalize when new feedback pathways were added. It produced alternating, stepping-like motion with the walker suspended in the air, presented at the 2021 NIH BRAIN Initiative Investigators Meeting \citep{bolenTorqueValuesArtificialKnee2021June15-17}, but it never produced walking: no hand-set combination of parameters bore the robot's weight, and the feedback pathways were never exercised under load. Those outcomes directly motivated the verification-first approach of the remainder of this chapter.

\section{What the AnimatLab Attempt Exposed}\label{sec:sns_gaps}

The air-stepping result and the hand-tuning failure together identified what the program lacked: capabilities that no single existing tool provided, and that hand tuning could not substitute for. Four gaps stood between the walker of Section~\ref{sec:prior_walker} and a controller development program for the BPA-actuated robot.

First, no environment coupled neural control to the models this program works on. The two-level controller existed only in AnimatLab projects; it had never been coupled to a full human-scale musculoskeletal model such as the Gait2392 model, nor to the robot-body model with its BPA muscle routing, and AnimatLab carried neither. MuJoCo, SNS-Toolbox, and the conversion tools existed as separate pieces with no assembled path between them \citep{todorov_mujoco_2012, nourse_sns_2023}.

Second, the contact-feedback walking circuits of the lineage were locked inside one engine. The Li walker and the Deng controller, with their foot-contact feedback, existed only as AnimatLab projects, and AnimatLab is no longer maintained. Circuits that exist in one unsupported engine cannot be checked against a second implementation, carried to a maintained physics engine, or run toward real-time hardware.

Third, there was no graphical circuit construction on maintained simulation tools. The circuits of this lineage are block diagrams, and building them by hand in code is slow and hard to inspect; a block-programming environment such as Simulink is the natural medium, and Simscape Multibody additionally accepts CAD-derived bodies. The predecessor SNS block libraries of this laboratory's own lineage \citep{hunt_development_2017, young_analyzing_2019} did not carry the elements the robot requires, and no current library did.

Fourth, no tool chain connected CAD-derived mechanics, BPA force models, and neural controllers in one loop. The robot exists as a SolidWorks assembly with BPAs whose forces are characterized in Chapter~\ref{ch:muscles}; connecting that geometry, those actuators, and a neural controller required interfaces that had to be built and verified.

These gaps also explain why the tools span several simulation engines rather than one. The engines differ in what they do best: MuJoCo runs fast scripted closed-loop simulation but is limited to rigid tendons, Simulink and Simscape Multibody provide graphical circuit construction and CAD import, and the MuJoCo--SNS-Toolbox combination supports real-time control and computer-in-the-loop operation for the hardware program. Porting the same controller across engines, and checking each port against its source, is how this chapter makes the circuits portable objects rather than artifacts of one simulator. The sections below state how far each gap is closed, and the verification tests state exactly what the assembled tools do and do not show.

\section{A Neuromechanical Simulation Toolchain}\label{sec:preliminary_sim_methods}

This section documents the tools and what each contributes: the existing tools they stand on, the new tools built from them, and the methodology by which the controllers are tuned and checked. Figure~\ref{fig:pipeline} diagrams the flow, from model conversion through closed-loop simulation to the embedded hardware stack of Chapter~\ref{ch:overall}. Quantitative results follow in Sections~\ref{sec:preliminary_sim_results} and~\ref{sec:spiking_results}.

\begin{figure}[htbp]
\centering
\includegraphics[width=\textwidth]{pipeline.pdf}
\caption[Simulation-to-robot pipeline]{The simulation-to-robot pipeline in three phases. Model conversion carries the OpenSim robot-body model through MyoConverter to a MuJoCo MJCF model with activation-state muscle actuators. Closed-loop simulation couples that model, with actuator forces drawn from the BPA force model of Chapter~\ref{ch:methods} inside the loop, to an SNS-Toolbox network through per-muscle activation commands and Ia, Ib, and II afferent plus heel and toe contact feedback. Hardware deployment, the future-work stage of Chapter~\ref{ch:overall}, moves the same loop onto Teensy-class microcontrollers for the per-joint valve, pressure, and contact loops and a Jetson Orin for network evaluation, computer-in-the-loop before embedded execution on the treadmill biped.}
\label{fig:pipeline}
\end{figure}

\subsection{Existing Tools and Data Resources}\label{sec:sns_existing_tools}

None of the infrastructure components is invented here: OpenSim supplies the musculoskeletal modeling side, the robot's source model, and the moment-arm and torque standards the conversions are audited against \citep{delp_opensim_2007, seth_opensim_2018}; MuJoCo supplies the physics engine \citep{todorov_mujoco_2012}; MyoConverter performs the base OpenSim-to-MuJoCo conversion \citep{caggiano_myosuite_2022, wang_myosim_2022}; SNS-Toolbox builds and integrates the spinal networks \citep{nourse_sns_2023}; AnimatLab hosts the controller lineage's archived reference models; Simulink and Simscape Multibody supply the controls-engineering environment; and SCONE provides the reference point for mature predictive neuromuscular simulation \citep{geijtenbeek_scone_2019}, whose published balance and gait demonstrations were reproduced headless as a practice baseline and whose balance tutorial anchors the standing stage of the training curriculum (Section~\ref{sec:sns_curriculum}). The sensory channels that close the loop, and the acceptance phenomena against which the controllers are validated, trace to the curated records of the Sensory Afferent Database of Chapter~\ref{ch:database} (Section~\ref{sec:sadb_use}).

\subsection{New Tools}\label{sec:sns_new_tools}

What is new is the combination, built as four tool families: the full-lower-body pipeline, the SNS-Simscape block library, the controller and body ports, and the verification instruments. Table~\ref{tab:sns_capability} states what each adds, the evidence that establishes it, and its completion status, distinguishing complete tools from partial imports, provisional parameters, and proposed entries; the subsections that follow record what the table cannot carry, each construction and its relevant limits.

\begin{table}[htbp]
\caption{New tools: what each adds, its verification evidence (Sections~\ref{sec:preliminary_sim_results}, \ref{sec:spiking_results}), and completion status. Partial imports assemble and run but carry a documented defect; provisional parameters are placeholder values not used for quantitative claims; proposed entries are specified but not built.}
\label{tab:sns_capability}
\small
\begin{tabularx}{\textwidth}{@{}>{\raggedright\arraybackslash}p{0.28\textwidth}>{\raggedright\arraybackslash}p{0.24\textwidth}>{\raggedright\arraybackslash}X@{}}
\toprule
\textbf{Tool (status)} & \textbf{Adds} & \textbf{Evidence} \\
\midrule
OpenSim-to-MuJoCo pipeline with BPA force interface (complete) & robot model with audited routing in a maintained engine; BPA forces inside the loop & 92 actuators audited; dynamic sign check; 410-neuron network at a \qty{2}{\milli\second} step \\
SNS-Simscape block library (complete; CAD parameters provisional) & graphical circuit construction; diameter-specific BPA force blocks; imported bodies & reflex demonstrations and spiking twins \\
Controller and body ports (complete) & Deng, W2L, and Li circuits across NumPy, Simulink, MuJoCo, Simscape & Deng period \qty{1.938}{\second} vs.\ \qty{1.94}{\second} reference; bridge readback bit-exact \\
SolidWorks knee into Simscape Multibody (partial import) & CAD-derived geometry under SNS control & hinge mates arrive welded; re-mate fix defined \\
\texttt{gait2392\_simbody} into Simscape Multibody (partial import) & full-body mechanism for Simulink-side control & 1,446-block mechanism compiles and simulates; muscles re-enter as blocks \\
Reorganized bilateral AnimatLab project (assembled, not yet oscillating) & bilateral rhythm generators and biarticular pathways & opens, exports, runs headless; latches under the drives tested \\
Cross-engine harness and gates (complete) & trust across engines & W2L rhythm core: period and burst counts match exactly \\
Spiking-neuron mirrors (complete) & coding-assumption test for every family & calibration residuals and gate tables (Appendix~\ref{app:neuron_equations}) \\
Shinohara, Di Russo, Shevtsova replications (proposed) & next replication series & circuit-rule translations with per-edge citations, not yet built \\
\bottomrule
\end{tabularx}
\end{table}

The SNS block library stands in a direct lineage, and the comparison matters. The SNS formalism descends from the block-based Simulink libraries of this laboratory's own research line, in which \citet{hunt_development_2017} trained the hindlimb-walking neural controller and with which \citet{young_analyzing_2019} built and analyzed a full-muscle rat hindlimb model; those libraries drive Hill-type biological muscles on animal-scale models. The library built here is a derivative, not a competitor: it keeps the same neuron and synapse equations so that circuits transfer, and it adds what the predecessors do not carry, Ia and Ib afferent encoders as blocks, diameter-specific BPA force elements implementing the fitted force surfaces of Chapter~\ref{ch:methods}, coupling to CAD-derived and OpenSim-derived bodies in Simscape Multibody, spiking additions, and regeneration from the tuned production network. A circuit assembled in the predecessor libraries controls a biologically realistic muscle model; a circuit assembled here can command the robot's actual actuators on its imported geometry.

\subsection{AnimatLab Body and Controller Organization}

The environment remains in the program because the controller lineage is archived in it: the reference models replicated in Section~\ref{sec:preliminary_sim_results} live as AnimatLab projects, and new circuit organizations are prototyped in the same format (Figures~\ref{fig:animatlab_model_architecture}--\ref{fig:animatlab_working_gui}, with the neuron-level motifs in Figure~\ref{fig:animatlab_controller_motifs}). The working model was instrumented for a \qty{10}{\second} phase-1 run, reported in Section~\ref{sec:preliminary_sim_results} and distinct from the earlier walker study of Section~\ref{sec:prior_walker}, and has since been reorganized into a complete bilateral network: the right-side subsystems mirrored from the left, inhibitory commissural pathways added between the two rhythm generators, and four biarticular muscle pathways added per leg with human maximum isometric forces from the Gait2392 data set. The reorganized project opens, exports, and runs headless, but holds a static latch rather than oscillating under the tonic drives tested, so conductance tuning is the active step and no locomotor claim is made beyond verified assembly.

\begin{figure}[p]
\centering
\includegraphics[width=0.64\textwidth,height=0.30\textheight,keepaspectratio]{OurWalker.png}
\par\medskip
\includegraphics[width=\textwidth,height=0.31\textheight,keepaspectratio]{2layerImprove.png}
\caption[AnimatLab walker and controller hierarchy]{AnimatLab model and architecture illustrations. Top: the segmented walker and labeled hip, knee, ankle, and toe joints. Bottom: the left/right subsystem organization and an expanded view of the left-side rhythm-generator, hip, and knee--ankle modules. These architectural views describe model organization and are not snapshots of the phase-1 trajectory.}
\label{fig:animatlab_model_architecture}
\end{figure}

\begin{figure}[p]
\centering
\begin{minipage}[t]{0.47\textwidth}\centering
\textbf{(A) Rhythm generator}\par\smallskip
\includegraphics[width=\linewidth,height=0.20\textheight,keepaspectratio]{RGnetwork.PNG}
\end{minipage}\hfill
\begin{minipage}[t]{0.47\textwidth}\centering
\textbf{(B) Pattern formation}\par\smallskip
\includegraphics[width=\linewidth,height=0.20\textheight,keepaspectratio]{PFnetwork.PNG}
\end{minipage}
\par\medskip
\textbf{(C) Sensorimotor organization}\par\smallskip
\includegraphics[width=0.76\textwidth,height=0.39\textheight,keepaspectratio]{sensoryMotor_Network.PNG}
\caption[AnimatLab controller reference schematics]{Reference schematics of \textbf{(A)} rhythm generation, \textbf{(B)} pattern formation, and \textbf{(C)} sensorimotor organization. Ext and Flx identify antagonist extensor and flexor pathways; IN denotes interneurons, MN motor neurons, and Ia and Ib muscle afferent pathways. Open triangles and filled circles distinguish excitatory and inhibitory synaptic connections. These schematics illustrate controller architecture and are not readouts from the phase-1 run.}
\label{fig:animatlab_controller_motifs}
\end{figure}

\begin{figure}[p]
\centering
\includegraphics[width=\textwidth,height=0.31\textheight,keepaspectratio]{animatlab_body_gui.png}
\par\medskip
\includegraphics[width=\textwidth,height=0.31\textheight,keepaspectratio]{animatlab_network_gui.png}
\caption[AnimatLab working-project interface]{The saved AnimatLab working project opened in the application: the biomechanical body at initialization (top) and the expanded left-side controller subsystem (bottom). The workspace tree and property inspector expose the relationship between the body, sensory pathways, and neural modules. These interface captures were acquired without running a new simulation.}
\label{fig:animatlab_working_gui}
\end{figure}

\subsection{MuJoCo Conversion and Synthetic BPA Knee}

The first MuJoCo tool moves the robot's musculoskeletal model from OpenSim into MuJoCo \citep{delp_opensim_2007, todorov_mujoco_2012}; the second makes the converted model's actuators behave as the BPAs of Chapter~\ref{ch:methods} rather than as generic muscles. MyoConverter performs the base conversion \citep{caggiano_myosuite_2022, wang_myosim_2022}; because an automatic conversion of a model with custom muscle routing introduces specific defects, among them paths not mirrored between the legs and joint axes of reversed sign, the conversion is followed by repair scripts that correct each defect and re-audit the model against the OpenSim source, and the pipeline runs in a recorded software environment so it can be re-run and reproduced. The force interface implements the quasi-static BPA force model of Chapter~\ref{ch:methods}, with the correction-factor force balance of Chapter~\ref{ch:torque}, inside the loop, reading every actuator's length and extension rate each step and returning its force.

The pair is complete and verified. The converted model carries all 92 muscle-tendon actuators with attachment geometry audited against the source, and the joint sign conventions were verified dynamically, each actuator activated alone, a check that must be dynamic because these attachment points move with joint configuration. An SNS-Toolbox network couples to MuJoCo at a \qty{2}{\milli\second} neural step, matching the physics timestep, with muscle force feedback entering as Ib afferent input; on this coupling a 410-neuron two-level spinal network has been assembled, and its rhythm generator has been verified in isolation to oscillate with the two legs in antiphase under tonic drive. The loop's sensory channels, Ia length and velocity, Ib force, group II length, and heel and toe contact marking the stance--swing boundary, follow from the curated records of the SADb (Chapter~\ref{ch:database}; Tables~\ref{tab:sadb_afferents} and~\ref{tab:sadb_pathways}) rather than from ad hoc literature searches at wiring time. Full-body closed-loop co-simulation runs under a simulated support harness that mirrors the laboratory's overhead rig; Section~\ref{sec:sns_controller_comparison} states exactly what walks and what does not. The pair's fundamental demonstration is a synthetic knee, a minimal SNS with reciprocal inhibition driving antagonistic flexor and extensor BPAs, which establishes the controller--actuator--joint interface and nothing more (Figure~\ref{fig:mujoco_model_views}).

\begin{figure}[p]
\centering
\begin{minipage}[t]{0.48\textwidth}\centering
\textbf{(A) Converted robot}\par\smallskip
\includegraphics[width=\linewidth,height=0.53\textheight,keepaspectratio]{mujoco_converted_robot.png}
\end{minipage}\hfill
\begin{minipage}[t]{0.48\textwidth}\centering
\textbf{(B) Synthetic BPA knee}\par\smallskip
\includegraphics[width=\linewidth,height=0.53\textheight,keepaspectratio]{mujoco_synthetic_knee.png}
\end{minipage}
\caption[MuJoCo model representations]{Native MuJoCo renders of \textbf{(A)} the converted Gait2392 robot at its saved reference pose, with muscle paths visible, and \textbf{(B)} the independent synthetic antagonist-BPA knee at an illustrative static angle of \ang{-35}. Cyan and gold identify its named flexor and extensor tendon routes; transparency exposes the routes through the view. The models are not shown to the same scale. These views document model construction, not SNS control of the converted robot or validated torque predictions.}
\label{fig:mujoco_model_views}
\end{figure}

\subsection{Simulink SNS Library and Reflex Topology}

The Simulink tool is a block library, built so that synthetic nervous systems can also be assembled and inspected in a standard controls-engineering environment. Seven core blocks, a graded neuron, a graded chemical synapse, a spiking leaky integrate-and-fire neuron, Ia and Ib afferent encoders, muscle activation dynamics, and a generic BPA force block, implement the same equations as the SNS-Toolbox networks; diameter-specific BPA force elements (\qtylist{10;20;40}{\mm}) implementing the fitted force surfaces of Chapter~\ref{ch:methods} and a Hill-type muscle element complete the library, which runs in base Simulink without additional toolboxes (Figure~\ref{fig:simulink_sns_implementation}).

Its fundamental demonstrations are the two canonical reflex pathways at the scale where their mechanisms are visible (Figure~\ref{fig:simulink_reflex_topology}), Ia stretch feedback with reciprocal inhibition and Ib force feedback with autogenic inhibition in a reduced-order knee with antagonist actuation, plus a beverage-holding arm loaded by a pour; their values are reported in Section~\ref{sec:sns_simulink_demos} beside their spiking twins. To couple the controller to detailed mechanical CAD, the SolidWorks assembly imports into Simscape Multibody, and the import carries one documented defect that bounds what the saved runs can claim: the exporter drops the SolidWorks hinge mates as unsupported, so the knee degree of freedom arrives welded, with the defined fix, re-mating each hinge as a concentric and coincident pair before re-export, not yet applied; the demonstration model also retains provisional mechanical parameters and a placeholder generic BPA force block, so the saved runs are not used to establish quantitative knee-regulation performance.

\begin{figure}[p]
\centering
\includegraphics[width=0.94\textwidth,height=0.42\textheight,keepaspectratio]{sns_library_overview.pdf}
\par\medskip
\includegraphics[width=0.94\textwidth]{simulink_neuron_detail.pdf}
\caption[Simulink SNS library blocks and neuron implementation]{Core SNS-Simscape blocks (top) and graded-neuron implementation (bottom). Summed currents drive membrane integration; threshold, inverse-slope gain, and saturation bound the output. Voltage feeds both current calculations. The generic BPA block remains a placeholder. These views do not demonstrate quantitative knee regulation or complete CAD coupling.}
\label{fig:simulink_sns_implementation}
\end{figure}

\begin{figure}[p]
\centering
\includegraphics[width=\textwidth]{KneeReflex_circuit.pdf}
\caption[Reduced-order knee-reflex circuit topology]{The reduced-order knee-reflex model drawn in the SNS library's marker language: triangular synaptic terminations mark excitatory inputs and filled round terminations inhibitory inputs. Ia stretch feedback excites the stretched muscle's motor-neuron pool and inhibits the antagonist pool (reciprocal inhibition); Ib force feedback inhibits the pool of the muscle that generated the force (autogenic inhibition). The motor-neuron pools drive the activation and BPA force blocks that act on the knee joint. The corresponding Simulink model assembles the library blocks one-for-one and simulates standalone. Mechanical parameters and the BPA force block remain provisional.}
\label{fig:simulink_reflex_topology}
\end{figure}

\subsection{Controller and Body Ports into Simulink and Simscape Multibody}\label{sec:sns_ports}

Making the controllers and bodies portable is the third tool family. The Deng two-layer CPG was rebuilt as separate editable Simulink files with the published parameter tables carried over \citep{deng_neuromechanical_2019, nourse_sns_2023}; one cell-level choice proved load-bearing for portability, a burst-termination property of the rhythm-generator cells whose membrane mechanism is recorded in Appendix~\ref{app:neuron_equations}. In the other direction, a script turns the tuned production network into an editable Simulink model built from the library blocks, one per neuron and one per synapse, regenerating whenever the tuning changes. Two body-import routes bring real geometry into Simscape Multibody: the SolidWorks knee via URDF and Multibody-Link XML, and the OpenSim \texttt{gait2392\_simbody} model via MuJoCo and URDF, with the massless muscle path-point slides preserved as links so that actuators attach at the same physical points as in the source; a URDF carries no muscles, so actuation re-enters as SNS library blocks driving the imported joints. A MuJoCo--Simulink bridge block set closes the loop from the Simulink side, muscle activation commands in and joint and actuator sensors out.

\subsection{Cross-Platform Verification of the Contact-Driven Rhythm Core}\label{sec:crossengine_methods}

The verification harness, a fixed script suite that re-runs and checks the models, is a tool in its own right, and its fundamental demonstration is replication in two forms: the same rhythm core rebuilt in two engines must reproduce the same rhythm, and the lineage's archived models must reproduce their published behavior in their original engine. The harness is defined around the contact-driven rhythm core of the Walker-2-Layer (W2L) bipedal architecture, four rhythm-generator populations, one extensor and one flexor population per leg plus graded pattern-formation and sensory cells, driven by heel-contact mechanosensory feedback \citep{deng_neuromechanical_2019}. The same core was implemented twice from a common connectivity specification, once as a NumPy reference through the SNS-Toolbox backend and once as native Simulink blocks, and both were driven with an identical protocol scored on burst period, left--right antiphase, burst counts, and full-trace correlation. The archived AnimatLab reference models were additionally re-run headless from their standalone simulators, so that engine-level and network-level agreement could be separated.

\subsection{Spiking-Neuron Mirrors of the Simulation Models}\label{sec:spiking_methods}

The final tool tests an assumption rather than adding a simulation capability. The controllers documented above compute with graded (non-spiking) neurons; to test that the conclusions drawn from them are not artifacts of graded coding, spiking mirrors of all three simulation families were built under one design rule: the motor-neuron layer stays non-spiking, because the musculoskeletal model consumes a continuous activation command, while the afferent encoders, rhythm-generator populations, and interneurons become leaky integrate-and-fire neurons operating at real membrane-potential ranges \citep{szczecinski_perspective_2023, nourse_sns_2023}. The MuJoCo mirror reproduces the production spinal network population for population, is selected at build time by a single environment-variable switch, and leaves the default build path untouched, proven by re-running its defaults gate and baseline evaluation to bit-exact agreement; because SNS-Toolbox 1.5.2 ships no spiking bursting-neuron class, its rhythm-generator populations are adapting integrate-and-fire cells. The mirrors' synaptic weights are set by a calibration script, not by hand, calibrated by mean conductance onto the motor neurons and by firing-rate transfer between spiking cells, with the full calibration tabulated in Appendix~\ref{app:neuron_equations}. In the Simulink family, three blocks were added, a \texttt{SpikingNeuron}, an event-wired \texttt{SpikingSynapse}, and a \texttt{HybridSpikingSynapse} carrying the same conductance dynamics onto a non-spiking postsynaptic cell, and spiking twins of the demonstrations replace only the sensory interneuron layer. In the AnimatLab family, a conversion script rewrites spiking copies of the runnable standalone exports, leaving every original untouched. Each mirror is verified against its non-spiking twin by a fixed gate set, topology (identical connection multiset), rhythm within \qty{20}{\percent} of the analog period, a basin gate under \qty{1}{\percent} multiplicative jitter on every scalar parameter, and regression to the committed artifacts, and carries mapped parameters only, with no new tuning, so its results establish what transfers structurally and what requires re-tuning.

\subsection{Tuning, Robustness, and Reference Data}\label{sec:sns_tuning_methods}

Above the physiologically set values, the free gains of each network are tuned by systematic search rather than by hand: a sequential model-based optimizer (a tree-structured Parzen estimator) proposes gain sets, each proposal is scored by a full simulation, and each stage's winner seeds the next. No hand-tuned parameter enters a production configuration; the failure mode of the prior walker was precisely hand tuning that left no record of sensitivity. Robustness is enforced during tuning rather than inspected afterward: finalist configurations must retain their rhythm under \qty{1}{\percent} multiplicative jitter on every scalar parameter, survive stepping of the tonic drive about the tuned value, and, for standing configurations, absorb push perturbations without falling. The requirement exists because the tuned networks sit near bifurcations of their dynamics, as the spiking results of Section~\ref{sec:spiking_results} make explicit. The reference data for tuning and evaluation are inverse-kinematics cycles from many human subjects of different sizes and speeds rather than a single subject: a 43-cycle locomotion reference library, detailed in Section~\ref{sec:gait_library_results}. Circuit structure is treated as data under the same methodology: every network's connection multiset is compared programmatically against its intended topology, and the publication schematics are generated from the compiled network with each drawn edge asserted against the compiled connection list.

\section{Verification: Tool Demonstrations and Experiments}\label{sec:preliminary_sim_results}

This section reports what the tools do, each subsection stating the expected behavior, the test or comparison, the observed result, and its limit. The BPA-actuated walker that demonstrates the full MuJoCo chain end to end is consolidated in Section~\ref{sec:sns_controller_comparison}.

\subsection{AnimatLab Phase-1 Recordings}

The phase-1 run of Section~\ref{sec:preliminary_sim_methods} is the AnimatLab tool's fundamental demonstration: it exercises the instrumented model and the recording chain end to end. Expected: synchronized, finite recordings over \qty{10}{\second}. Observed: joint, neural, contact-sensory, and body-motion channels recorded at a uniform \qty{0.2}{\milli\second} interval, with bilateral hip, knee, and ankle angles shown in Figure~\ref{fig:animatlab_phase1_preliminary}. Limit: the contact-sensory channels record membrane voltage rather than contact force, the model included an externally applied horizontal force during the first second, and the run demonstrates signal acquisition and alternating activity, not validated locomotor performance.

\begin{figure}[htbp]
\centering
\includegraphics[width=\textwidth,height=0.60\textheight,keepaspectratio]{animatlab_phase1_joint_angles.pdf}
\caption[AnimatLab phase-1 hip, knee, and ankle recordings]{AnimatLab phase-1 joint motion: \textbf{(A)} hip, \textbf{(B)} knee, and \textbf{(C)} ankle angles in the model's recorded joint conventions. Blue solid and orange dashed traces denote left and right sides. All 50,001 samples within the configured 0--\qty{10}{\second} chart window are retained without filtering. The model includes an external horizontal force during the first second. These traces document joint motion in an instrumented run; they do not establish stable walking or effective ground-contact gating.}
\label{fig:animatlab_phase1_preliminary}
\end{figure}

\subsection{Simulink Library Demonstrations and Cross-Program Ports}\label{sec:sns_simulink_demos}

The SNS-Simscape library demonstrates the canonical reflex pathways at the scale where their mechanisms are visible: the knee-reflex joint settles at a mean of \ang{43.5} over its last \qty{2}{\second}; the stretch reflex limits the beverage arm's sag to \ang{2.00} against \ang{7.85} with the reflex disabled; and a third demonstration closes the loop onto the actuator model itself, a burst-terminating two-population CPG driving the antagonist BPA pair of the one-degree-of-freedom knee through the diameter-specific force blocks. The committed values are tabulated beside their spiking twins in Section~\ref{sec:spiking_results} (Table~\ref{tab:spike_knee}).

The controller ports verify at the same standard. The Deng files, driven by a single current pulse, oscillate at a \qty{1.938}{\second} period against the \qty{1.94}{\second} NumPy reference, with the antagonist populations anticorrelated at $r = -0.86$; the exported full-network model reproduces its NumPy reference to \qty{4.2e-06}{\milli\volt} across all 410 neurons over the early transient, beyond which the network's chaos makes any two integrators diverge in phase. The MuJoCo--Simulink bridge passes single-muscle, two-rate, and sensor-readback tests against an independent Python integration, the readback bit-exact over the scored window; on the bridge, a single-muscle Ia/II/Ib reflex loop on the converted model's vastus medialis bounds a postural collapse from transient muscle forces past \qty{40}{\kilo\newton} down to \qty{469}{\newton}, and the full exported network ran as a Simulink block driving all 92 muscle actuators through a \qty{5}{\second} closed-loop run, isolating a real property: under tonic drive alone the Simulink realization settles to the tonic fixed point while the NumPy reference oscillates (Section~\ref{sec:disc_spiking}). The \texttt{gait2392\_simbody} model arrives in Simscape Multibody as a 1,446-block, 149-joint mechanism that compiles and simulates, with the muscles re-entering as SNS library blocks driving the imported joints.

\subsection{Cross-Engine Verification of the Contact-Driven Rhythm Core}\label{sec:crossengine_results}

Expected: the same rhythm core, implemented twice from a common specification, reproduces the same rhythm. Observed: between the NumPy reference and the native Simulink reimplementation of the W2L rhythm core, both legs' burst periods matched exactly at \qty{1.000}{\second}, the left--right antiphase correlation and all ten expected bursts reproduced the reference, and the full voltage traces cross-correlate at $r \geq 0.99996$ with an RMSE below 0.111~mV (Figure~\ref{fig:w2l_crossengine}); the residual lag is consistent with the one-step delay of the discrete Simulink solver. The same protocol also runs the chapter's first deletion experiment, and it surfaced an architectural fact rather than a defect: with the heel-contact trains removed, the full network holds a static latch under tonic drive with kickoff alone, in every implementation, so phasic contact feedback is part of the drive architecture of the W2L design, not an optional perturbation, which is why the pathway-removal experiments of Section~\ref{sec:walker_variant_results} treat the contact channels as a pathway class rather than as a boundary condition.

\begin{figure}[htbp]
\centering
\includegraphics[width=\textwidth,height=0.55\textheight,keepaspectratio]{sns_w2l_traces.png}
\caption[W2L rhythm-core traces across engines]{Left and right rhythm-generator extensor population voltages of the W2L contact-driven rhythm core under the shared verification protocol, from the NumPy reference (SNS-Toolbox backend) and the native Simulink reimplementation. Scored metrics (burst period, left--right antiphase correlation, and burst counts) agree exactly, and the full traces cross-correlate at $r \geq 0.99996$ at a $-\qty{2}{\milli\second}$ lag.}
\label{fig:w2l_crossengine}
\end{figure}

\subsection{Headless Validation of the AnimatLab Reference Models}\label{sec:animatlab_validation}

Expected: the lineage's archived models reproduce their published rhythms and gaits in their original engine. Observed: all eleven targetable rhythm and gait metrics pass their acceptance bands. The bilateral contact-driven model's rhythm generators oscillate at \qty{0.450}{} and \qty{0.467}{\second} with a half-cycle left--right lag; the modern W2L model's joints oscillate at \qty{2.25}{\Hz} with clean extensor/flexor antiphase ($r = -0.996$) carried by the graded pathway rather than by flexor spiking; and the Li-model rerun \citep{li_bipedal_2017} walks at \qty{0.637}{\meter\per\second} with step periods near \qty{1.3}{\second}, reproducing the published behavior of the original model in its original engine. Limit, stated plainly: the recent connectome-level circuit organizations of \citet{shevtsova_modeling_2016}, \citet{dirusso_investigating_2023}, and \citet{shinohara_mechanisms_2025} are specified in this program as circuit-rule translations, connection lists carrying per-edge citations, from which running models would be built; they have not been built or verified, and no replication claim is made for them here.

\subsection{Gait-Validation Sweep against the Locomotion Reference Library}\label{sec:gait_library_results}

The library comprises 32 running trials from eight subjects at four speeds, ten walking references (seven at prescribed speeds from \qty{0.5}{\metre\per\second} to \qty{2.0}{\metre\per\second}, plus three self-selected), and the predicted-walking cycle of the Falisse predictive-simulation dataset.

Expected: a network tuned to one walking reference transfers, at least partially, to diverse human reference locomotion. Observed: the four production walking configurations, spanning the program's two controller lineages, were scored against the 43-cycle library with the production pipeline's cycle-shape kinematic score, in which 0 denotes perfect agreement and $-25$ is the fixed no-cycle fallback (Table~\ref{tab:gait_sweep}). The ranking agrees across evaluations, with the W2L-layout variant closest and the six-synergy variant farthest, but every configuration's library scores are one to two orders of magnitude more negative than its tuning-reference score, a distance that is itself the finding: it quantifies how far a network tuned to a single walking reference remains from the library's diverse cycles. All four configurations produce right-leg-only or asymmetric gait in these evaluations (Section~\ref{sec:sns_controller_comparison}).

\begin{table}[htbp]
\caption{Gait-validation sweep of the four production walking configurations against the 43-cycle locomotion reference library. The score is the production pipeline's cycle-shape kinematic score (0 denotes perfect agreement); the tuning-reference score is each configuration's evaluation on the tuning pipeline's walking reference, and the library mean averages the same score over all 43 references (among genuine gait cycles, higher scores indicate closer agreement).}
\label{tab:gait_sweep}
\begin{tabular}{@{}lrrl@{}}
\toprule
 & \textbf{Tuning-ref.} & \textbf{Library mean} & \textbf{Most negative reference} \\
 & \textbf{score} & \textbf{score (43 refs)} & \textbf{(score)} \\
\midrule
\texttt{s3k}        & $-186.4$ & $-8776$ & \texttt{Case\_40\_motion} ($-8841$) \\
\texttt{s3kpruned}  & $-176.1$ & $-6924$ & \texttt{Case\_40\_motion} ($-6986$) \\
\texttt{syn6}       & $-200.4$ & $-9961$ & \texttt{Case\_40\_motion} ($-10040$) \\
\texttt{w2lvar}     & $-144.5$ & $-5808$ & \texttt{subject08\_Run\_20002} ($-5896$) \\
\bottomrule
\end{tabular}
\end{table}

\subsection{Stimulus and Deletion Experiments on the Walker Networks}\label{sec:walker_variant_results}

The locomotor physiology that this program's circuits imitate was established with deletion and stimulation interventions, among them proprioceptive-deletion degradation of locomotor patterns in mice \citep{akay_degradation_2014}, regulation of locomotor activity by proprioceptive input \citep{lam_role_2002}, and tonic stimulation that facilitates involuntary air-stepping in humans \citep{selionov_tonic_2009}; the assembled tools make the simulator-side analogues inexpensive, removal of an entire pathway class from a fully tuned walker, re-evaluated on the tuning reference, and stepping of the tonic drive about its tuned value. The systematic stimulation series at the pattern-formation and motoneuron levels is defined as future verification work (Chapter~\ref{ch:futurework}) and has not yet been run.

The pathway-removal sweep (Table~\ref{tab:walker_ablations}) removes individual pathway classes from each configuration and re-evaluates the tuning-reference score. In the production configuration, removing the Ia afferent pathway freezes the model outright, with no gait cycles detected and the fixed $-25$ fallback returned. Removing the Renshaw cells, the group-II afferents, or all afferent classes at once degrades the score to $-267.3$, $-255.2$, and $-267.8$, respectively, against the full model's $-186.4$, while removing the Ib, contact, or interleg commissural channels individually shifts the score by smaller amounts toward zero. The six-synergy variant behaves differently: with all afferents removed it scores $-190.5$ against its full model's $-200.4$ and still walks bilaterally, with ten cycles on each leg, a knee-flexion minimum of \ang{-71.2}, and a \qty{0.684}{\second} period. The W2L-layout variant depends on both the Ia pathway (removed: $-273.5$) and the weak rhythm-generator mutual excitation (removed: $-279.9$) against its full model's $-144.5$, while removing its Renshaw cells leaves a bilateral gait at $-142.4$.

\begin{table}[htbp]
\caption{Pathway-removal sweep on the tuning-reference evaluation. The score is the production pipeline's cycle-shape kinematic score (0 denotes perfect agreement; $-25$ is the fixed no-cycle fallback, a failed evaluation that scores above every genuine gait in this table). The full-model rows repeat each configuration's unmodified production evaluation.}
\label{tab:walker_ablations}
\begin{tabular}{@{}llr@{}}
\toprule
\textbf{Configuration} & \textbf{Removed pathway} & \textbf{Score} \\
\midrule
\texttt{s3k} (full model) & none & $-186.4$ \\
 & Ia afferents & $-25$ (no cycles detected) \\
 & all afferent classes & $-267.8$ \\
 & Renshaw cells & $-267.3$ \\
 & group-II afferents & $-255.2$ \\
 & Ib afferents & $-180.3$ \\
 & contact channels & $-183.9$ \\
 & interleg commissural pathways & $-174.9$ \\
\addlinespace
\texttt{syn6} (full model) & none & $-200.4$ \\
 & all afferent classes & $-190.5$ \\
\addlinespace
\texttt{w2lvar} (full model) & none & $-144.5$ \\
 & Ia afferents & $-273.5$ \\
 & weak rhythm-generator mutual excitation & $-279.9$ \\
 & Renshaw cells & $-142.4$ \\
\bottomrule
\end{tabular}
\end{table}

Read against the reference physiology, deleting all afferent classes from the production configuration degrades its gait without abolishing it, the qualitative outcome of the deafferentation experiments \citep{akay_degradation_2014}, matched most closely by the six-synergy variant, while the Ia-only removal freezing that network marks its tuning as built around its length feedback, a stronger dependence than the animal experiments show. On the stimulus side, the production configuration never fell in 15 stand-and-push evaluations and survives drive perturbations of $\pm 0.01$ and $-0.02$, freezing only at $+0.02$, with its pruned derivative surviving the same ladder, while the six-synergy and W2L-layout variants fall under every perturbation tested. The sweep identifies the production configuration and its pruned derivative as the robust pair of the family and locates the fragility of the two variants under drive perturbation.

\section{Spiking-Neuron Mirror Results}\label{sec:spiking_results}

The spiking mirrors of Section~\ref{sec:spiking_methods} were evaluated against their non-spiking twins in the order in which the families were introduced: the MuJoCo spinal network, the Simulink demonstrations, and the AnimatLab walker family. All numbers come from the gated verification effort whose artifacts and reproduction commands are catalogued in Appendix~\ref{app:neuron_equations}.

\subsection{MuJoCo Spinal-Network Mirror}

The non-spiking production baselines were first re-established on the recorded software environment: all three production configurations reproduce their documented evaluation scores exactly, and after the mirror's selector was added, the default build gate and that same baseline score ($-161.57$) reproduce digit for digit. The topology gate passes exactly: under the default, fully tuned, and joint-layer gain configurations, the spiking mirror's directed connection multiset is identical to the non-spiking build's, with identical input lists and motor-neuron maps.

Under constant drive (network alone, \qty{20}{\second}), the rhythm gate splits. The spiking mirror self-sustains clean bilateral antiphase, with rhythm-generator extensor and flexor cells firing at \qty{4.9}{\hertz} and \qty{3.2}{\hertz} per leg, but its envelope period is \qty{0.621}{\second} against the analog build's \qty{0.399}{\second}, a \qty{55.6}{\percent} lengthening that misses the gate's \qty{20}{\percent} tolerance (Figure~\ref{fig:spike_rhythm}). The basin gate passes in full: the rhythm persists under \qty{1}{\percent} jitter on every scalar parameter in all 12 seeded trials (final-window swing minimum \qty{4.58}{\mV} against the \qty{5.13}{\mV} baseline, margin $0.89\times$). In the deafferented suspended-air preparation the mirror air-steps rhythmically with full-amplitude joint swings (knee range \ang{106.0} against the twin's \ang{117.9}), but its cadence is roughly ten times faster than the twin's \qty{0.47}{\hertz} regime, so the air-regime gate fails on the cadence clause alone.

\begin{figure}[htbp]
\centering
\includegraphics[width=\textwidth]{spiking_mirror_rhythm.pdf}
% ALT TEXT: see spiking_mirror_rhythm_alt.txt beside the figure file.
\caption[Constant-drive rhythm, non-spiking versus spiking]{Constant-drive rhythm of the right rhythm-generator under network-only simulation. \textbf{(A)} The non-spiking network's extensor-minus-flexor readout voltage alternates irregularly about a \qty{4.8}{\milli\volt} swing (larger only in the startup transient), with an average period of \qty{0.399}{\second}. \textbf{(B)} The spiking mirror's raw extensor (solid orange) and flexor (dashed indigo) cells fire in alternating bursts whose envelope period is \qty{0.621}{\second}, \qty{55.6}{\percent} longer than the analog build.}
\label{fig:spike_rhythm}
\end{figure}

One untuned ground-walking evaluation was run as a stretch attempt: the spiking mirror at the non-spiking winner's mapped parameters, with no re-tuning. It scores worse overall (kinematic score $-236.35$ against $-161.57$), as expected for an untuned mapped starting point, but it produces a genuine bilateral gait (eight right-leg and nine left-leg cycles, a cycle-window knee flexion minimum of \ang{-79.8}, no non-finite states, and the model stays up), whereas the non-spiking winner's left leg is frozen in every recorded run of its study (Figure~\ref{fig:spike_ground}, Table~\ref{tab:spike_ground}). This is a single run with no statistics, reported as a lead, not as an established gait result.

\begin{table}[htbp]
\caption{Untuned spiking mirror versus its non-spiking twin at the same winner parameters, ground evaluation (16 s). The cycle knee minimum and stance fractions are computed over detected gait cycles; the kinematic score is the cycle-shape error metric of the production pipeline (0 denotes perfect agreement; higher scores indicate closer agreement). The spiking mirror carries mapped parameters only, with zero re-tuning.}
\label{tab:spike_ground}
\begin{tabular}{@{}lcc@{}}
\toprule
\textbf{Metric} & \textbf{Non-spiking winner} & \textbf{Spiking mirror (untuned)} \\
\midrule
Kinematic score                & $-161.57$          & $-236.35$ \\
Gait cycles, right / left      & 9 / 0 (left frozen) & 8 / 9 (bilateral) \\
Cycle knee minimum             & \ang{-64.3}         & \ang{-79.8} \\
Contact duty, right / left     & 0.637 / 0.985       & 0.799 / 0.856 \\
Maximum pelvis tilt            & \ang{22.1}          & \ang{19.2} \\
Non-finite states / fell       & no / up             & no / up \\
\bottomrule
\end{tabular}
\end{table}

\begin{figure}[htbp]
\centering
\includegraphics[width=\textwidth]{spiking_mirror_ground_gait.pdf}
% ALT TEXT: see spiking_mirror_ground_gait_alt.txt beside the figure file.
\caption[Ground gait, non-spiking winner versus untuned spiking mirror]{Ground-walking joint angles of the converted MuJoCo robot at the non-spiking winner's parameters. \textbf{(A, C)} The non-spiking winner steps with the right leg only; the left knee moves through a single slow excursion without cycling. \textbf{(B, D)} The untuned spiking mirror cycles both legs with deep knee flexion. Solid indigo traces denote the left leg, and dashed orange traces denote the right leg.}
\label{fig:spike_ground}
\end{figure}

\subsection{Simulink Spiking Demonstration Twins}

The three new library blocks were verified at the unit level before any demonstration was rebuilt: single-spike and pulse-train conductance transients match the exact solutions to the tabulated residuals, and a hybrid two-neuron circuit matches a verbatim re-implementation of the SNS-Toolbox update equations spike for spike; with the blocks present, every committed demonstration and units test reproduces its recorded value exactly.

At the matched operating point, the spiking twin of the knee-reflex demonstration is indistinguishable from the baseline in its transient and within one degree in its settle behavior (Table~\ref{tab:spike_knee}, Figure~\ref{fig:spike_knee}): both read \ang{15.79} at \qty{0.05}{\second}, and the twin settles \ang{0.9} higher, because the rate-coded pathway adds the synaptic filter time constant that graded synapses do not have. The beverage-holding twin likewise holds its reflex modulation within \ang{0.02} of the baseline, with the reflex-on sag at \ang{2.02} against the baseline's \ang{2.00} and the reflex-off sag identical.

\begin{table}[htbp]
\caption{Knee-reflex demonstration, committed baseline versus spiking twin at the matched operating point (5 s runs; settle statistics over the last \qty{2}{\second}).}
\label{tab:spike_knee}
\begin{tabular}{@{}lcc@{}}
\toprule
\textbf{Metric} & \textbf{Baseline} & \textbf{Spiking twin} \\
\midrule
Knee angle at \qty{0.05}{\second}   & \ang{15.79} & \ang{15.79} \\
Settle mean (last \qty{2}{\second}) & \ang{43.5}  & \ang{44.4} \\
Settle range                        & \ang{41.6}--\ang{44.9} & \ang{41.9}--\ang{46.7} \\
Extensor / flexor activation (window means) & 0.53 / 0.53 & 0.58 / 0.56 \\
Antagonist alternation              & \qty{0.5}{\hertz}--\qty{2}{\hertz} & \qty{7.4}{\hertz} (dominant peak) \\
Ia interneuron firing at settle     & graded, subthreshold & \qty{95}{\hertz}; Ib \qty{0}{\hertz} \\
\bottomrule
\end{tabular}
\end{table}


\begin{figure}[htbp]
\centering
\includegraphics[width=\textwidth]{simulink_knee_spiking_twin.pdf}
% ALT TEXT: see simulink_knee_spiking_twin_alt.txt beside the figure file.
\caption[Simulink knee-reflex twin, baseline versus spiking]{Knee angle of the Simulink knee-reflex demonstration: committed non-spiking baseline (solid indigo) versus spiking twin (dashed orange) at the matched operating point. \textbf{(A)} Full \qty{5}{\second} run. \textbf{(B)} Settle window (last \qty{2}{\second}); the spiking twin's alternation band is slightly wider.}
\label{fig:spike_knee}
\end{figure}

A representative rhythm-generator-to-motor-pool subnetwork demonstrates the hybrid motif end to end, spiking populations driving a non-spiking motor-neuron pool whose graded voltage tracks the rhythm generator's rate envelope; its burst alternation is partial, reported as is, and the full 410-population Simulink spiking conversion is future work for the same reason.

\subsection{AnimatLab Spiking Conversions}

The AnimatLab family produces the clearest negative result of the spiking comparison. The graded baselines were re-established first: the two-layer walker's rhythm generator oscillates at \qty{0.443}{\second}, the bilateral variant at \qty{0.462}{\second} with the two legs folded near antiphase ($r = 0.869$), and the already-spiking 2023-lineage model, the family's spiking reference, bursts at \qty{1.305}{\second}, reproducing its documented period. No spiking copy of the three graded models preserves the rhythm (Table~\ref{tab:spike_animatlab}): the two-layer conversions die at every threshold and decay tried; the bilateral variant latches with one side's populations saturated at \qty{442}{\hertz}--\qty{450}{\hertz}; and the reorganized project's copy converts its latched baseline into strictly synchronous \qty{11.8}{\hertz} co-bursting with no alternation. The thresholds-only ablation, which keeps the graded synapses, preserves near-baseline joint oscillations, isolating the synapse conversion as the change that breaks the rhythm: per-spike conductances at the graded amplitude either starve the reciprocal loop of inhibitory charge or saturate it, whereas the graded synapses read the membrane potential continuously and keep oscillating through subthreshold voltages. A spiking version of this family therefore needs its synaptic amplitudes and decays re-tuned as a rate code, which is how the 2023 spiking reference was built in the first place. The conversions are catalogued variant by variant in Appendix~\ref{app:neuron_equations} (Table~\ref{tab:app_spk_animatlab_full}), together with the full gate table of the MuJoCo mirror.

\begin{table}[htbp]
\caption{AnimatLab spiking conversions versus their graded baselines (5.1 s headless runs; joint ranges in degrees where the rhythm survives). The thresholds-only ablation converts the neurons but keeps the graded synapses.}
\label{tab:spike_animatlab}
\small
\begin{tabular}{@{}lcc@{}}
\toprule
\textbf{Variant} & \textbf{Baseline (graded)} & \textbf{Spiking copy} \\
\midrule
Two-layer walker     & \qty{0.443}{\second} rhythm;   & dead (2 spikes; \\
                     & hip $-15$ to $+23$, knee $-1$ to $+60$ & hip pinned near $+6$) \\
Bilateral variant    & \qty{0.462}{\second}; antiphase $-0.42$ cycle & latched; one side at \qty{442}{\hertz}--\qty{450}{\hertz} \\
Reorganized project  & latched (tonic drives) & synchronous \qty{11.8}{\hertz} co-bursting; \\
(standalone export)  &                        & same-leg envelopes $r = 1.000$ at lag 0 \\
Thresholds only      & \multirow{2}{*}{hip $-15$ to $+23$, knee $-1$ to $+60$} & oscillates: hip $-15.1$ to $+23.2$, \\
(graded synapses)    &                        & knee $-0.8$ to $+47.9$ \\
\bottomrule
\end{tabular}
\end{table}

\section{Delivered Capabilities, Limits, and the Program They Enable}\label{sec:sns_discussion}

This closing section states what the spiking comparison means for the program's neural-coding assumption, records the training curriculum behind the tuned configurations, and consolidates the BPA-actuated robot, the standing of the robot models, and the walking status, before handing the program off to Chapter~\ref{ch:futurework}.

\subsection{What Spiking Changes and Does Not}\label{sec:disc_spiking}

What spiking does not change is the behavior of the controlled model, provided the design rule and the calibration hold: because the motor-neuron layer stays non-spiking and the spikes onto it are calibrated by mean conductance, the continuous activation interface that the actuator and torque models of this dissertation attach to is unchanged by construction. The Simulink twins make the point quantitatively, with the knee-reflex settle band moving by \ang{0.9} and the beverage-holding reflex modulation by \ang{0.02}. The actuator-level and joint-level conclusions of Chapters~\ref{ch:methods}--\ref{ch:results} therefore do not depend on the neural coding chosen above the motor-neuron layer.

What spiking changes is timing and robustness. The mirror's self-sustained rhythm runs \qty{55.6}{\percent} slower at the same constant drive, and its air-stepping cadence runs roughly ten times faster, deviations in opposite directions, which a single scaling error could not produce; the constant-drive period traces to two identifiable knobs, the threshold-increment magnitude and the adaptation time constant, so a period-matched spiking controller is a short tuning effort away rather than a structural failure. Robustness moves in the mirror's favor: the spiking rhythm-generator populations passed the perturbation basin gate in full where the analog network's rhythm sits close to a bifurcation, and for the hardware-transfer program of Chapter~\ref{ch:futurework}, where the network will be re-implemented in finite-precision arithmetic, a wide basin is the safer substrate. What spiking does not tolerate is a naive parameter copy at the synapse level: the AnimatLab conversions show that converting graded amplitudes to per-spike increments of the same magnitude breaks the reciprocal rhythm-generator loop, while converting only the thresholds changes nothing; synaptic strength in a spiking network is a rate-code quantity that must be derived from the presynaptic firing rate and the conductance decay, exactly as the calibration script derives it. The one place where the untuned mirror outperforms its tuned twin, a bilateral ground gait where the winner's left leg is frozen, is a single run with no statistics, claimed only as a lead for the closed-loop CPG comparison program of Chapter~\ref{ch:futurework}.

\subsection{The Five-Stage Training Curriculum and Its Scoring}\label{sec:sns_curriculum}

The tuned configurations of this chapter were not tuned all at once. The training proceeds through five stages (Figure~\ref{fig:sns_curriculum}), each adding one capability and seeded with the previous stage's winner so that only the newly activated pathway gains are searched: deafferented air stepping first, afferented air stepping second, standing balance third, ground walking fourth, and ground walking with the contact-driven pathway variant fifth. Standing before walking is the design rule: gait tuning begins only after a configuration holds posture, and the failures the curriculum surfaces early are the cheap ones. The scoring is stage-appropriate and gated: rhythm and swing depth for the air stages, with the objective rhythm-gated so a frozen configuration cannot win by holding a pose; posture maintenance under push for the balance stage; and the cycle-shape kinematic score for the walking stages. The walking score is rhythm-gated in the same way for a subtler reason: a run with no detectable gait cycles returns a fixed sentinel ($-25$) rather than a score, because that sentinel sits numerically above every genuine gait score, and reading it as a score would let a frozen model out-score a walker. Scores are never read alone: finalists must pass the perturbation and basin requirements of Section~\ref{sec:sns_tuning_methods}, because a score achieved on a narrow basin is a property of the solver rather than of the circuit.

\begin{figure}[htbp]
\centering
\includegraphics[width=\textwidth]{sns_curriculum.pdf}
\caption[Five-stage training curriculum and stage scores]{The five-stage training curriculum. Each stage adds one capability and is seeded with the previous stage's winner, so only the newly activated pathway gains are searched; the numbered rows distinguish the preparations and their objectives. The air-stepping stages are suspended, and the reported ground evaluations retain the support harness. Scores are stage-appropriate and rhythm-gated where rhythm is the capability under test, and finalists at every stage must pass the perturbation and basin gates before the score is read.}
\label{fig:sns_curriculum}
\end{figure}

\subsection{The BPA-Actuated Robot and the Walking Status}\label{sec:sns_controller_comparison}

The actuator model travels with the controller: inside the MuJoCo loop, actuator force comes from the quasi-static BPA model of Chapter~\ref{ch:muscles} with the correction-factor force balance of Chapter~\ref{ch:torque}, and in Simulink the library's diameter-specific force blocks implement the same fitted surfaces, so every closed-loop demonstration in this chapter draws its forces from one validated source. What the integration does not yet carry is actuator dynamics: the pressure--force dynamics under pulse-modulated valve control \citep{elzein_pressure_2025} enter the program as future work (Chapter~\ref{ch:futurework}).

The end-to-end demonstration of the MuJoCo chain, conversion pipeline plus force interface plus network builder, is the BPA-actuated robot. The walker geometry of Section~\ref{sec:prior_walker} was exported to MuJoCo and driven by antagonistic BPA actuator pairs under the split contact-driven network, with actuator forces computed by the BPA model. In the air-stepping preparation, with the harness carrying 70\% of the \qty{411}{\N} body weight, the walker produced clean knee-flexion excursions at a \qty{1.542}{\second} period with a knee flexion range of \ang{63.5}, peak BPA forces of \qty{880}{\N} (left) and \qty{1372}{\N} (right) at a peak pressure of \qty{372}{\kPa} (60\% of the \qty{620}{\kPa} ceiling), and 7 and 5 heel-contact episodes on the left and right: the complete chain from the characterized BPA force model to a multi-joint walker stepping under synthetic neural control.

The same evaluation series reported its failures with equal weight. On the ground, the W2L gate model under its untuned default gains stands (full-support stance) but does not walk: heel duty is zero and forward displacement is $+\qty{0.001}{\meter}$ over the run. With the left and right drives separated and the feedback gains re-tuned, the same robot alternates in antiphase (inter-leg phase 0.487) and progresses \qty{1.72}{\meter} over a \qty{30}{\second} window (\qty{0.066}{\metre\per\second}) under the harness protocol, with the neural drive weaned to \qty{60}{\percent} of its bootstrap value and the result robust under perturbation (12 of 12 basin re-runs); below \qty{60}{\percent} drive the gait dies, the harness itself is not weaned, and foot duty remains below the strict walking gate, so the result is supported progression rather than validated walking. The MuJoCo re-implementation of the Li model produced no stance episodes in any of seven drop-protocol variants. Together with the latch result of Section~\ref{sec:crossengine_results}, these negatives localize the open problem precisely: the contact-driven rhythm architecture itself transfers across engines intact, but translating it to unsupported ground locomotion requires the balance and weight-bearing layers defined in Chapter~\ref{ch:futurework}, and those layers are not present in these models.

The walking status deserves a precise statement, because the chapter's records support exactly this and no more. Each tuned configuration was brought through the curriculum's air and balance stages, and on the ground all walking evaluations run under the support harness, where the configurations produce gross walking-like progression with right-leg-only or asymmetric gait patterns. A subsequent drive-weaning study re-tuned the feedback gains while forcing the hand-drive down a ladder of scales: its finalist improves the default build's evaluation score from $-161.57$ to $-145.07$ with a bilateral gait of ten contact cycles per leg (contact duty 0.71 and 0.65), survives a closed-loop basin check in eight of eight perturbed evaluations, and, with its feedback held fixed, keeps bilateral stepping at every drive scale down to zero hand-drive, at which the drive pathway is absent from the network entirely, whereas the un-retuned configuration collapses below \qty{85}{\percent} of full drive. The gait remains in-place alternation under the support rig, without forward progression. Unsupported, harness-free walking has not been demonstrated anywhere in this program's records, and whether any configuration walks without the harness is an open question; the balance and weight-bearing layers that would answer it are the first item of the simulation program defined in Chapter~\ref{ch:futurework}.

On the converted robot-body model, the four production walking configurations span the controller lineages, and scored with the same cycle-shape kinematic score against the 43-reference library, the family-level fact is that the robot-walker-lineage configuration, not the human-derived ones, agrees most closely with the library, while the six-synergy configuration, the most strongly human-derived in its pattern formation, agrees least closely; closing that gap is a tuning-program question rather than a tool question. On the physical-model side, the SolidWorks knee assembly is imported into Simscape Multibody with its degree-of-freedom defect documented and its fix defined, and the imported \texttt{gait2392\_simbody} mechanism awaits its muscle blocks; neither physical model yet carries a tuned controller.

What the tools are for is the program they define next: the balance and weight-bearing layers that the walking status below calls for, the stimulation and deletion series at the pattern-formation and motoneuron levels, and the pressure--force dynamics under valve control, all defined in Chapter~\ref{ch:futurework}. Every tool is reusable by a researcher who did not build it, and the loop's speed is what makes the deletion and stimulus experiments cheap enough to repeat.
